\documentclass[12pt,a4paper,oneside]{book}

%==============================================================================
% ENCODING AND FONTS
%==============================================================================
\usepackage[utf8]{inputenc}
\usepackage[T1]{fontenc}
\usepackage{newpxtext}
\usepackage{newpxmath}
\usepackage[vietnamese,english]{babel}
\usepackage[scaled=0.95]{inconsolata}

%==============================================================================
% PAGE LAYOUT AND TYPOGRAPHY
%==============================================================================
\usepackage[a4paper, margin=1in, bindingoffset=0.5cm]{geometry}
\usepackage[final,protrusion=true,expansion=true]{microtype}
\usepackage{setspace}
\onehalfspacing
\usepackage{emptypage}

%==============================================================================
% MATH
%==============================================================================
\usepackage{mathtools}
\usepackage{amssymb}

%==============================================================================
% GRAPHICS AND FIGURES
%==============================================================================
\usepackage{graphicx}
\usepackage[font=small,labelfont=bf,format=hang]{caption}
\usepackage{subcaption}

%==============================================================================
% TABLES
%==============================================================================
\usepackage{booktabs}
\usepackage{array}
\usepackage{multirow}

%==============================================================================
% LISTS
%==============================================================================
\usepackage{enumitem}

%==============================================================================
% COLORS
%==============================================================================
\usepackage{xcolor}
\definecolor{chaptercolor}{RGB}{0, 70, 140}
\definecolor{linkblue}{RGB}{0,51,153}

%==============================================================================
% HEADER/FOOTER
%==============================================================================
\usepackage{fancyhdr}
\pagestyle{fancy}
\fancyhf{}
\fancyhead[L]{\small\nouppercase{\leftmark}}
\fancyhead[R]{\small\thepage}
\makeatletter
\let\headrule\@headrule
\makeatother
\fancypagestyle{plain}{
    \fancyhf{}
    \fancyfoot[C]{\small\thepage}
}

%==============================================================================
% CHAPTER TITLE STYLING
%==============================================================================
\usepackage{titlesec}
\titleformat{\chapter}[display]
  {\normalfont\huge\bfseries\color{chaptercolor}}
  {\chaptertitlename\ \thechapter}{20pt}{\Huge}
\titlespacing*{\chapter}{0pt}{-20pt}{40pt}

%==============================================================================
% ALGORITHMS
%==============================================================================
\usepackage{algorithm}
\usepackage{algpseudocode}

%==============================================================================
% THEOREMS
%==============================================================================
\usepackage{amsthm}
\theoremstyle{plain}
\newtheorem{theorem}{Theorem}[chapter]
\newtheorem{definition}[theorem]{Definition}
\newtheorem{example}[theorem]{Example}

%==============================================================================
% UNITS AND NUMBERS
%==============================================================================
\usepackage{siunitx}
\sisetup{detect-all}

%==============================================================================
% QUOTATIONS
%==============================================================================
\usepackage[autostyle]{csquotes}

%==============================================================================
% CITATIONS (BibTeX / natbib)
%==============================================================================
\usepackage{natbib}
\setcitestyle{square,numbers,sort&compress}

%==============================================================================
% APPENDICES
%==============================================================================
\usepackage{appendix}

%==============================================================================
% HYPERLINKS
%==============================================================================
\usepackage{bookmark}
\usepackage{hyperref}
\hypersetup{
    colorlinks=true,
    linkcolor=chaptercolor,
    citecolor=linkblue,
    urlcolor=linkblue,
    pdfauthor={Author Name},
    pdftitle={Thesis Title},
    pdfsubject={Information Systems},
    pdfkeywords={data platform, knowledge graph, intelligent analytics},
    bookmarks=true,
    bookmarksnumbered=true,
    bookmarksopen=true,
}

%==============================================================================
% SMART CROSS-REFERENCES
%==============================================================================
\usepackage{cleveref}

%==============================================================================
% TABLE OF CONTENTS / LISTING PACKAGES
%==============================================================================
\usepackage{tocbibind}
\usepackage{minitoc}

%==============================================================================
% CUSTOM COMMANDS
%==============================================================================
\newcommand{\R}{\mathbb{R}}
\newcommand{\N}{\mathbb{N}}
\newcommand{\Z}{\mathbb{Z}}
\DeclareMathOperator*{\argmax}{arg\,max}
\DeclareMathOperator*{\argmin}{arg\,min}
\DeclarePairedDelimiter{\abs}{\lvert}{\rvert}
\DeclarePairedDelimiter{\norm}{\lVert}{\rVert}

%==============================================================================
% CUSTOM TITLE PAGE (VNU / UET FORMAT)
%==============================================================================
\usepackage{ifthen}
\newboolean{isDoctoral}
\setboolean{isDoctoral}{false} % Set to true for Doctoral, false for Master's

%==============================================================================
\begin{document}

% ============================================================
% FRONT MATTER — Roman numerals, no chapter numbering
% ============================================================
\frontmatter
\dominitoc

% --- TITLE PAGE ---
\begin{titlepage}
\begin{center}

% University Name
\vspace*{1.5cm}
{\Large\textbf{VIETNAM NATIONAL UNIVERSITY HANOI}}\\[0.3cm]
{\Large\textbf{UNIVERSITY OF ENGINEERING AND TECHNOLOGY}}\\[2cm]

% Author name
{\LARGE\bfseries AUTHOR NAME}\\[2cm]

% Thesis title
{\LARGE\bfseries\color{chaptercolor}
THESIS TITLE\\
Building a Modern Data Platform\\
Based on Data Fabric Architecture}\\[2cm]

% Degree / Major
{\large Major: Information Systems}\\[1.5cm]

% Supervisor
{\large Supervisor:}\\
{\large Associate Professor NAME}\\[2cm]

% Year / Location
{\large Ha Noi -- YEAR}

\vfill

\end{center}
\end{titlepage}

% --- BLANK PAGE ---
\clearpage\thispagestyle{empty}\null\clearpage

% --- ACKNOWLEDGEMENTS ---
\chapter*{Acknowledgements}
\addcontentsline{toc}{chapter}{Acknowledgements}

% PLACEHOLDER: Replace with your acknowledgements
I would like to express my sincere gratitude to the professors at the University of Engineering and Technology, Vietnam National University Hanoi, for their dedicated guidance and teaching throughout my studies and research at the university. I would like to extend my deep appreciation to my supervisor and all the esteemed faculty members of the Information Systems Department. They have provided me with valuable guidance and direction in completing this thesis.

Despite my best efforts to complete the thesis to the best of my ability, there are still some shortcomings. I sincerely hope to receive feedback and suggestions from my professors, colleagues, and friends to improve the thesis.

Thank you very much!

\vspace{1cm}
\hfill Ha Noi, Month YEAR

\clearpage

% --- AUTHORSHIP DECLARATION ---
\chapter*{Authorship}
\addcontentsline{toc}{chapter}{Authorship}

I hereby declare that all the results reported in this thesis are the outcome of my own work under the guidance of my supervisor.

All references to related studies have been clearly cited and listed in the reference section of the thesis. No part of this thesis has been copied from any other work without clear acknowledgment of the source. All experimental results reported in this thesis have been conducted and statistically analyzed based on the actual experimental data.

\vspace{2cm}
\hfill Author Name

\clearpage

% --- ABSTRACT (English) ---
\chapter*{Abstract}
\addcontentsline{toc}{chapter}{Abstract}

% PLACEHOLDER: Replace with your abstract (150-300 words)
In today's digital age, organizations are generating and collecting more data than ever before. This data comes from a variety of sources and is stored in different formats and locations. Managing this vast amount of data and making it accessible to the right people at the right time is becoming increasingly challenging. A modern data platform based on data fabric or knowledge graph architecture provides a unified, flexible, and scalable view of all enterprise data.

This thesis aims to design and implement a metadata knowledge graph-based data platform that can be easily integrated in modern platforms. This research also develops intelligent analytics features, which is an LLM module supporting querying and generating visualizations.

The main contributions of this thesis include: (1) introducing the definition of metadata knowledge graph, intelligent analytics and benefits of integrating it into an intelligent data platform; (2) designing the system architecture and technology stacks; (3) implementing and installing the platform; (4) evaluating the platform's effectiveness with a specific use case.

% PLACEHOLDER: Keywords
\noindent\textbf{Keywords:} Data Platform, Knowledge Graph, Intelligent Analytics, Data Fabric, Metadata Management

\clearpage

% --- TÓM TẮT (Vietnamese) ---
\chapter*{Tóm tắt}
\addcontentsline{toc}{chapter}{Tóm tắt}

% PLACEHOLDER: Replace with your Vietnamese abstract
Trong thời đại số ngày nay, các tổ chức đang tạo ra và thu thập lượng dữ liệu ngày càng lớn từ nhiều nguồn khác nhau. Việc quản lý lượng dữ liệu khổng lồ này và làm cho nó có thể truy cập được bởi đúng người vào đúng thời điểm đang trở nên thách thức hơn bao giờ hết. Một nền tảng dữ liệu hiện đại dựa trên kiến trúc data fabric hoặc knowledge graph cung cấp một cái nhìn thống nhất, linh hoạt và có thể mở rộng cho tất cả dữ liệu doanh nghiệp.

Luận văn này nhằm thiết kế và triển khai một nền tảng dữ liệu dựa trên knowledge graph và metadata, có thể dễ dàng tích hợp vào các nền tảng hiện đại. Nghiên cứu này cũng phát triển các tính năng phân tích thông minh, là mô-đun LLM hỗ trợ truy vấn và tạo trực quan hóa.

Các đóng góp chính của luận văn bao gồm: (1) giới thiệu định nghĩa về knowledge graph dựa trên metadata, phân tích thông minh và lợi ích của việc tích hợp vào nền tảng dữ liệu thông minh; (2) thiết kế kiến trúc hệ thống và các ngăn xếp công nghệ; (3) triển khai và cài đặt nền tảng; (4) đánh giá hiệu quả của nền tảng với một kịch bản cụ thể.

% PLACEHOLDER: Keywords Vietnamese
\noindent\textbf{Từ khóa:} Nền tảng dữ liệu, Knowledge Graph, Phân tích thông minh, Data Fabric, Quản lý Metadata

\clearpage

% --- TABLE OF CONTENTS ---
\tableofcontents
\clearpage

% --- LIST OF FIGURES ---
\listoffigures
\clearpage

% --- LIST OF TABLES ---
\listoftables
\clearpage

% --- LIST OF ABBREVIATIONS ---
\chapter*{List of Abbreviations}
\addcontentsline{toc}{chapter}{List of Abbreviations}
\begin{acronym}[Trino]
  \acro{ACID}{Atomicity, Consistency, Isolation, Durability}
  \acro{API}{Application Programming Interface}
  \acro{CDC}{Change Data Capture}
  \acro{DAG}{Directed Acyclic Graph}
  \acro{ETL}{Extract, Transform, Load}
  \acro{IoT}{Internet of Things}
  \acro{KG}{Knowledge Graph}
  \acro{LLM}{Large Language Model}
  \acro{MDM}{Master Data Management}
  \acro{ML}{Machine Learning}
  \acro{NoSQL}{Not Only SQL}
  \acro{PII}{Personally Identifiable Information}
  \acro{SQL}{Structured Query Language}
\end{acronym}
\clearpage

% ============================================================
% MAIN MATTER — Arabic numerals, chapter numbering
% ============================================================
\mainmatter

% --- INTRODUCTION (thesis-level, not a chapter) ---
\chapter*{Introduction}
\label{ch:introduction}
\addcontentsline{toc}{chapter}{Introduction}
% =============================================================================
% INTRODUCTION
% Chapter-level intro — NOT a numbered chapter
% =============================================================================
% COMPLETELY REWRITTEN — Old Knowledge Graph / Data Fabric content removed
% Project: A Context-Aware Framework for Streaming Data Quality Monitoring
% Date: April 24, 2026
% =============================================================================

\section{Motivation and Problem Statement}
\label{sec:intro-motivation}

Modern transportation systems — including urban bus fleets, taxi services, and
rail networks — generate continuous streams of Global Positioning System (GPS)
position records at update intervals ranging from seconds to minutes. These records
constitute the primary data source for real-time operations such as arrival-time
prediction, fleet management, and service scheduling. The quality of these records
directly determines the reliability of downstream applications; erroneous or
anomalous GPS data can propagate through the entire operational pipeline,
degrading predictive accuracy and compromising service quality.

GPS position records in public transit systems are susceptible to multiple classes of
data quality defects. Sensor noise produces implausible speed measurements;
GPS spoofing can inject fabricated vehicle positions separated by hundreds of meters
within seconds; transmission errors cause duplicate event records; and stale
position updates may propagate outdated location information to downstream
applications. These defects are not merely theoretical concerns — the NYC MTA Bus
system, for example, operates hundreds of vehicles updating at approximately
30-second intervals across the five boroughs of New York City, generating
millions of records per day that must be continuously monitored for quality.

Existing data quality frameworks were designed primarily for batch processing
environments. Great Expectations \citep{great_expectations} and Soda Core
\citep{soda-core} provide sophisticated constraint specification and validation
pipelines, but both operate in batch mode: data must be written to a warehouse
before quality checks can execute. In batch processing, anomalies are detected
hours after they occur, making these tools unsuitable for real-time transit
monitoring where immediate detection is operationally critical. Great Expectations
and Soda Core additionally lack cross-record validation capabilities that would
enable detection of GPS trajectory anomalies such as impossible speed jumps.

Stream DaQ \citep{streamdaq2025} represents the most closely related prior work,
providing a streaming-native data quality framework with real-time violation
detection. However, Stream DaQ was designed for generic data streams and lacks
GPS trajectory validation capabilities; it cannot detect physically implausible
vehicle speeds or coordinate jumps in transit GPS data. METER \citep{meter2024}
addresses concept drift in streaming machine learning models but does not implement
data quality validation for trajectory data. GTFS Validator \citep{gtfs_validator}
validates static GTFS schedule data but operates in batch mode and does not
support real-time GTFS-realtime feeds.

The gap we identify is the absence of a streaming data quality framework that
implements cross-record GPS trajectory validation on real-time GTFS-realtime
feeds. No surveyed framework — among Great Expectations, Soda Core, Stream DaQ,
METER, Ada-Context \citep{adacontext2025}, Deequ \citep{deequ2018}, GTFS Validator,
or any of 14 others surveyed in our literature review — implements Haversine-based
GPS speed and jump detection on streaming vehicle positions with context-aware
threshold adaptation. This gap is operationally significant: erroneous GPS data
in public transit directly degrades arrival-time predictions and fleet management
decisions that affect thousands of passengers daily.

We present \textbf{StreamDQ}: a context-aware framework for streaming data quality
monitoring that validates GPS trajectory quality on real GTFS-realtime feeds in
real time. StreamDQ is not a production-grade system; it is a research and
education platform designed to demonstrate that streaming GPS validation is feasible,
measurable, and more effective than batch tools for real-time transit data quality.

\section{Research Objectives}
\label{sec:intro-objectives}

We formulate six research questions that address the core claims of this work.
Each question is stated with a measurable hypothesis and a target threshold;
all thresholds are labeled \texttt{[TIER-2\ ESTIMATED]} because they require
benchmark validation.

\textbf{RQ1.} Do hierarchical context-aware thresholds (L0--L4) improve detection F1
over static global thresholds? We hypothesize a target improvement of
$\Delta\text{F1} \geq 5$ percentage points for events falling into L0 context
cells (fine-grained hour--zone--weekday), assessed via Wilcoxon signed-rank test
with Holm--Bonferroni correction ($\alpha_{\text{adj}} = 0.0083$).
\texttt{[TIER-2\ ESTIMATED]}

\textbf{RQ2.} Does the L4 global fallback reduce threshold specificity
detectably compared to L0 context-specific thresholds? We assess whether the
false-negative rate at L4 is higher than at L0 by more than 10\%,
measured by Friedman's test across context cells.
\texttt{[TIER-2\ ESTIMATED]}

\textbf{RQ3.} Does Trajectory Quality Scoring (TQS) correlate with a
downstream quality signal? We measure Pearson correlation $\rho > 0.70$ between
TQS composite scores and independent downstream prediction error. Critically,
we redesign RQ3 to avoid circularity: TQS is computed from raw event properties
alone (not from rule violation counts), and correlation is measured against
ETA prediction error on NYC MTA Bus, not against injection rate.
\texttt{[TIER-2\ ESTIMATED]}

\textbf{RQ4.} Does the cross-record (CRS) rule layer provide incremental F1
beyond what syntactic and semantic rules achieve alone? We target
$\Delta\text{F1} \geq 5$ percentage points from adding CRS rules to the SYN+SEM
baseline, measured via Wilcoxon signed-rank test.
\texttt{[TIER-2\ ESTIMATED]}

\textbf{RQ5.} Do CRS rules achieve precision $> 0.70$ on synthetically injected
GPS anomalies in the NYC MTA Bus GTFS-realtime stream? We assess both CRS001
(GPS speed bounds) and CRS002 (GPS jump detection) using Wilson score confidence
intervals at 95\% confidence level.
\texttt{[TIER-2\ ESTIMATED]}

\textbf{RQ6.} Does ML-augmented threshold calibration (Bayesian Optimization)
improve F1 over rule-only thresholds? This is a measured contribution:
positive, negative, and mixed results are all scientifically valuable and will
be reported as such. Phase 3 implementation is mandatory.
\texttt{[TIER-2\ ESTIMATED]}

\section{Contributions}
\label{sec:intro-contributions}

We implement four primary contributions for the StreamDQ research platform.

\textbf{C1: First streaming GPS trajectory validation on real GTFS-realtime
feeds.} We implement CRS001 (Haversine-based GPS speed bounds $[2.0,\ 100.0]$
km/h) and CRS002 (GPS jump detection $>400$ m in 30 s) as Java
\texttt{KeyedProcessFunction} operators in Apache Flink, evaluated on the live
NYC MTA Bus GTFS-realtime public feed. No prior framework implements cross-record
GPS validation on streaming transit data.
\texttt{[TIER-2\ ESTIMATED]}

\textbf{C2: Hierarchical context-aware threshold fallback (L0--L5).} We design
a five-level hierarchical threshold system that traverses from the finest-grained
context cell (L0: hour $\times$ zone $\times$ weekday/weekend, $\geq 100$
samples) to physics priors (L5), gracefully degrading through intermediate levels
(L1--L4). We validate that this fallback mechanism is absent from all
surveyed frameworks.
\texttt{[TIER-2\ ESTIMATED]}

\textbf{C3: Ground-truth evaluation methodology with explicit limitations.}
We implement a complete evaluation pipeline comprising synthetic anomaly injection,
per-entity ground-truth tracking in PostgreSQL, and precision/recall/F1 metrics
with 95\% bootstrap confidence intervals (1,000 iterations). Uniquely, we
explicitly label what cannot be measured: CRS003 recall is blocked by a known
evaluation bug (NG-4), and distributed Flink latency is not benchmarked.
\texttt{[TIER-1\ VERIFIED\ for\ methodology;\ TIER-2\ ESTIMATED\ for\ results]}

\textbf{C4: ML-augmented threshold calibration.} We design a three-component
ML pipeline — Bayesian Optimization (Priority 1, GO), Isolation Forest (Priority 2,
conditional), and XGBoost (Priority 3, conditional) — that calibrates context
thresholds without overriding rule-based violation detection. LSTM is explicitly
excluded (Priority 4, NO-GO) due to unconfirmed training data and high
redundancy with CRS002.
\texttt{[TIER-2\ ESTIMATED]}

\section{Thesis Structure}
\label{sec:intro-structure}

The remainder of this thesis is organized as follows.

\textbf{Chapter 1} presents the literature review and theoretical background,
covering streaming data quality definitions and frameworks, context-aware data
quality methods, GPS trajectory validation for transportation systems, and the
role of machine learning in data quality calibration.

\textbf{Chapter 2} describes the StreamDQ system architecture and algorithm
design, including the three-layer rule taxonomy (syntactic, semantic, and
cross-record), the hierarchical context-aware threshold system, the Trajectory
Quality Scoring mechanism, and the ML augmentation pipeline.

\textbf{Chapter 3} details the experimental evaluation methodology, including
datasets, synthetic anomaly injection, ground-truth tracking, statistical tests,
and the ablation study design. Result sections contain placeholder tables
labeled \texttt{[TIER-2\ ESTIMATED]} pending benchmark execution. A mandatory
limitations section explicitly acknowledges seven documented constraints.

The \textbf{Conclusion} summarizes the contributions and key findings, and
identifies five directions for future work, including fixing the CRS003
evaluation blocker, implementing the external context stub, and deploying a
distributed Flink cluster for full-scale benchmarks.

\clearpage

% --- CHAPTER 1 ---
\chapter{Fundamental Theories}
\label{ch:fundamentals}
% =============================================================================
% CHAPTER 1: FUNDAMENTAL THEORIES
% Chương 1: CƠ SỞ LÝ THUYẾT VÀ TỔNG QUAN
% =============================================================================

\section{Giới thiệu}
\label{sec:ch1_intro}

% PLACEHOLDER: Write your chapter introduction here
Chương này trình bày các cơ sở lý thuyết và tổng quan về các công nghệ được sử dụng trong luận văn. Nội dung bao gồm giới thiệu về Knowledge Graph, Data Fabric, Data Lakehouse, và các thành phần của nền tảng dữ liệu hiện đại.

\section{Knowledge Graph}
\label{sec:knowledge_graph}

\subsection{Khái niệm Knowledge Graph}
\label{sec:kg_concept}

% PLACEHOLDER: Define Knowledge Graph
Knowledge Graph (Đồ thị Tri thức) là một cách biểu diễn tri thức dưới dạng đồ thị, trong đó các nút (nodes) đại diện cho các thực thể và các cạnh (edges) đại diện cho các quan hệ giữa chúng~\citep{paulheim2017knowledge}.

\begin{definition}[Knowledge Graph]
Knowledge Graph là một cấu trúc dữ liệu dạng đồ thị có hướng bao gồm tập hợp các thực thể, các thuộc tính của thực thể, và các quan hệ giữa các thực thể.
\end{definition}

\subsection{Ứng dụng của Knowledge Graph}
\label{sec:kg_applications}

Knowledge Graph được ứng dụng rộng rãi trong nhiều lĩnh vực:

\begin{itemize}
    \item \textbf{Tìm kiếm thông minh:} Google Knowledge Graph, Bing Satori
    \item \textbf{Hệ thống khuyến nghị:} Amazon, Netflix sử dụng knowledge graph để đề xuất sản phẩm
    \item \textbf{Quản lý dữ liệu doanh nghiệp:} Metadata management và data lineage
    \item \textbf{Hỏi đáp tự động:} Chatbots và virtual assistants
\end{itemize}

\section{Data Fabric và Data Lakehouse}
\label{sec:data_fabric_lakehouse}

\subsection{Data Fabric Architecture}
\label{sec:data_fabric}

% PLACEHOLDER: Describe Data Fabric
Data Fabric là một kiến trúc dữ liệu tổng hợp, cung cấp khả năng truy cập dữ liệu thống nhất từ nhiều nguồn khác nhau thông qua các layer trừu tượng~\citep{van_der_vliet_2021_datafabric}.

\begin{figure}[tbp]
\centering
\fbox{\parbox{0.8\textwidth}{\centering
\vspace{2cm}
[Data Fabric Architecture Diagram]
\vspace{2cm}
}}
\caption{Kiến trúc Data Fabric tổng quan}
\label{fig:data_fabric}
\end{figure}

\subsection{Data Lakehouse}
\label{sec:data_lakehouse}

% PLACEHOLDER: Describe Data Lakehouse
Data Lakehouse là mô hình kiến trúc kết hợp ưu điểm của Data Lake và Data Warehouse truyền thống~\citep{armbrust2021lakehouse}.

\begin{itemize}
    \item \textbf{Scalability:} Lưu trữ dữ liệu lớn với chi phí thấp
    \item \textbf{Flexibility:} Hỗ trợ nhiều định dạng dữ liệu (structured, semi-structured, unstructured)
    \item \textbf{Performance:} Truy vấn nhanh với data skipping và caching
    \item \textbf{Governance:} Quản lý truy cập và lineage dữ liệu
\end{itemize}

\section{Các thành phần của Nền tảng Dữ liệu}
\label{sec:platform_components}

% PLACEHOLDER: Describe platform components
Một nền tảng dữ liệu hiện đại bao gồm các thành phần chính sau:

\begin{enumerate}
    \item \textbf{Data Source Layer:} Nguồn dữ liệu từ nhiều hệ thống operational (MySQL, PostgreSQL, MongoDB)
    \item \textbf{Data Ingestion Layer:} Thu nạp dữ liệu thời gian thực và batch (Apache Kafka, Debezium, Airbyte)
    \item \textbf{Storage Layer:} Lưu trữ dữ liệu thô (MinIO, HDFS, S3)
    \item \textbf{Transformation Layer:} Xử lý và biến đổi dữ liệu (Apache Spark, dbt)
    \item \textbf{Metadata Layer:} Quản lý metadata và knowledge graph (Apache Atlas, DataHub)
    \item \textbf{Analytics Layer:} Phân tích và trực quan hóa (Trino, Metabase, Superset)
\end{enumerate}

\section{Tổng kết chương}
\label{sec:ch1_summary}

% PLACEHOLDER: Write chapter summary
Chương này đã trình bày các cơ sở lý thuyết về Knowledge Graph, Data Fabric và Data Lakehouse. Các khái niệm và công nghệ được giới thiệu trong chương này sẽ là nền tảng cho việc thiết kế và triển khai hệ thống ở các chương tiếp theo.

\clearpage

% --- CHAPTER 2 ---
\chapter{Proposed System Architecture}
\label{ch:architecture}
% ============================================================
% CHAPTER 2: PROPOSED SYSTEM / ARCHITECTURE
% Trích xuất từ: 2 thesis mẫu tại thes_template/
% Cấu trúc: Related Work → Core Layers → Technology Architecture → Summary
% ============================================================

\chapter{Proposed System Architecture}
\label{ch:architecture}

% ============================================================
% SECTION 2.1: RELATED WORK AND BASE ARCHITECTURE
% ============================================================
\section{Related Work and Base Architecture}
\label{sec:related_work}

% ==========================================================
% SUBSECTION 2.1.1: Reference Architecture
% ==========================================================
\subsection{Reference Architecture}
\label{sec:reference_arch}

% PLACEHOLDER: Giới thiệu kiến trúc tham chiếu
% Tham khảo thesis 2023: vDFKM framework
% Tham khảo thesis 2026: mở rộng vDFKM

Our system architecture is referenced by a framework called [Reference Framework Name]~\citep{framework2023}. This is a standard platform based on [architecture type] applied in enterprise and organizational environments when it is necessary to build a centralized, secure, and fast data platform for in-depth analysis purposes.

\Cref{fig:reference_arch} shows the architecture of the reference framework.

\begin{figure}[tbp]
\centering
\fbox{\parbox{0.8\textwidth}{\centering\vspace{2cm}
[Insert reference architecture diagram here]\\
Architecture diagram showing layers: Data Source → Ingestion → Storage → Transformation → Metadata → Analytics
\vspace{2cm}}}
\caption{Reference framework architecture.}
\label{fig:reference_arch}
\end{figure}

% ==========================================================
% SUBSECTION 2.1.2: Proposed Architecture Overview
% ==========================================================
\subsection{Proposed High-Level Architecture}
\label{sec:proposed_arch}

% PLACEHOLDER: Mô tả kiến trúc đề xuất cao cấp
% 5 tầng chính + mô tả từng tầng

Based on the reference architecture, the proposed platform consists of five major layers that work together to support [goal]. At the base, the [Layer 1] leverages [technology] to [capability]. These changes are [processed/stored] into the platform and passed along to the next layer.

The [Layer 2] integrates [technologies] for [capability]. This layer enables both [capability 1] and [capability 2] access to raw and curated data.

Above this, the [Layer 3] utilizes [technologies] for [capability]. [Orchestration tool] orchestrates the entire pipeline, ensuring automated, reliable data workflows.

The semantic context and governance of the data are managed in the [Layer 4]. [Tool] plays a central role by [capability]. [Validation tool] enhances this layer with data validation and profiling.

At the top, the [Layer 5] enables end-users to interact with the data through [tools]. Integrated [technology] allows users to pose natural language questions, which are translated into executable queries.

\Cref{fig:proposed_arch} illustrates the proposed high-level architecture.

\begin{figure}[tbp]
\centering
\fbox{\parbox{0.9\textwidth}{\centering\vspace{2.5cm}
[Insert proposed architecture diagram here]\\
Show all 5 layers with arrows indicating data flow
\vspace{2.5cm}}}
\caption{Proposed high-level architecture.}
\label{fig:proposed_arch}
\end{figure}

% ============================================================
% SECTION 2.2: CORE LAYERS OF THE DATA PLATFORM
% ============================================================
\section{Core Layers of the Data Platform}
\label{sec:core_layers}

% ==========================================================
% 2.2.1 Data Source Layer
% ==========================================================
\subsection{Data Source Layer}
\label{sec:data_source_layer}

The Data Source Layer is where all the platform's data sources are centered. It can include a variety of different sources, ranging from relational databases to APIs, IoT sensors, web logs, or enterprise applications. These data sources can be structured, semi-structured, or unstructured.

\textbf{Structured data} is organized in a specific format (tables, columns), with predefined constraints, and easy to manage through traditional database management systems.

\textbf{Semi-structured data} has some organization but not enough to be considered structured. Formats such as XML, JSON, and HTML are common.

\textbf{Unstructured data} typically lacks any rules or organization and can be in various formats such as social media posts, image files, audio files, or unstructured text.

\Cref{tab:data_sources} summarizes the data sources used in this thesis.

\begin{table}[tbp]
\centering
\caption{Data sources used in the platform.}
\label{tab:data_sources}
\begin{tabular}{@{}lll@{}}
\toprule
\textbf{Source} & \textbf{Type} & \textbf{Format} \\
\midrule
Source A & Relational DB & Structured \\
Source B & API & Semi-structured \\
Source C & File Storage & Unstructured \\
\bottomrule
\end{tabular}
\end{table}

% ==========================================================
% 2.2.2 Data Ingestion Layer
% ==========================================================
\subsection{Data Ingestion Layer}
\label{sec:data_ingestion_layer}

The Data Ingestion Layer is the bridge between the Data Source Layer and downstream systems. This layer performs the process of importing, acquiring data from various sources for further analysis.

Data Ingestion tasks can be divided into two types:

\begin{description}
    \item[\textbf{Batch ingestion:}] Processes large volumes of data at once, typically on a scheduled basis. This method is useful when large amounts of data need to be ingested.
    \item[\textbf{Streaming ingestion:}] Processes data as it is generated, in real-time. This method is useful when data needs to be analyzed in real-time, such as in fraud detection or sensor monitoring.
\end{description}

Tools for data ingestion include open-source solutions like [Tool A], [Tool B], and commercial offerings from vendors like [Vendor A] and [Vendor B].

% ==========================================================
% 2.2.3 Data Storage and Lakehouse Layer
% ==========================================================
\subsection{Data Storage and Lakehouse Layer}
\label{sec:storage_layer}

The Data Storage and Lakehouse Layer is the foundational component of any modern data platform. This layer is responsible for storing and managing data, including raw, semi-structured, and processed datasets.

There are two popular types of storage architecture: \textbf{block storage} (used for high-performance transactional systems) and \textbf{object storage} (designed for scalable, distributed data access).

Object storage stores data in individual units at the same level. It does not use folders or subdirectories to organize data. Each object includes metadata, which provides information about the file and helps with its processing and usability.

\textbf{Data lake}, \textbf{data warehouse}, and \textbf{data lakehouse} are three different approaches to managing and processing large amounts of data:

\begin{itemize}
    \item \textbf{Data lake:} Storage repository for large amounts of raw data in any format.
    \item \textbf{Data warehouse:} Centralized repository optimized for querying and reporting.
    \item \textbf{Data lakehouse:} Hybrid architecture combining the best of both data lakes and data warehouses, enabling ACID transactions.
\end{itemize}

% ==========================================================
% 2.2.4 Data Transformation Layer
% ==========================================================
\subsection{Data Transformation Layer}
\label{sec:transformation_layer}

The Data Transformation Layer plays a critical role in the data stream by converting raw data into clean and structured formats suitable for analytics and downstream processing.

Transformation involves:

\begin{itemize}
    \item \textbf{Data cleaning:} Removing unwanted or irrelevant data, dropping unused columns, filtering duplicate values.
    \item \textbf{Data normalization:} Ensuring values are represented in consistent units and formats.
    \item \textbf{Data aggregation:} Calculating metrics such as averages, sums, and counts.
    \item \textbf{Schema alignment:} Mapping varying structures into unified formats.
\end{itemize}

Modern transformation workflows may also incorporate real-time streaming capabilities to enable continuous updates of data lakes.

% ==========================================================
% 2.2.5 Metadata and Knowledge Graph Layer
% ==========================================================
\subsection{Metadata and Knowledge Graph Layer}
\label{sec:metadata_layer}

Metadata and knowledge graph layer is one of the most important layers in data platforms by organizing and contextualizing metadata with meaningful links. This layer integrates with all the other layers to aggregate metadata throughout the system and visualize them into a knowledge graph.

This layer will recognize entities such as datasets, jobs, dashboards, and the relationships between them. Advanced properties of entities such as owners or tags will also need to be carefully configured.

From a practical standpoint, the Metadata and Knowledge Graph Layer supports advanced capabilities such as data discovery, impact analysis, and recommendation systems. It provides access to information about datasets, including schemas, data quality, update frequency, business definitions, ownership, and lineage.

\Cref{fig:kg_layer} shows the knowledge graph structure in the platform.

\begin{figure}[tbp]
\centering
\fbox{\parbox{0.7\textwidth}{\centering\vspace{2cm}
[Insert knowledge graph diagram here]\\
Nodes: datasets, jobs, dashboards; Edges: lineage, ownership, uses
\vspace{2cm}}}
\caption{Knowledge graph structure in the platform.}
\label{fig:kg_layer}
\end{figure}

% ==========================================================
% 2.2.6 Analytics and Visualization Layer
% ==========================================================
\subsection{Analytics and Visualization Layer}
\label{sec:analytics_layer}

This layer includes visualization tools that allow users to explore and comprehend their data visually. These technologies help users uncover trends, patterns, and anomalies in their data and efficiently communicate insights to stakeholders.

An important aspect in this layer supporting analytics tasks is the [LLM module], which helps translate user questions into queries and commands to generate visualizations. Instead of queries like ``SELECT SUM(...) FROM...'', users can simply ask the module to perform specific analytical tasks.

This approach makes it easier for people to get started and allows more people across the organization to access data. These approaches work best when paired with strong semantic storage that provides aggregated metrics.

% ==========================================================
% 2.2.7 Data Flow in the Platform
% ==========================================================
\subsection{Data Flow in the Platform}
\label{sec:data_flow}

\Cref{fig:data_flow} illustrates the complete data flow within the platform.

\begin{figure}[tbp]
\centering
\fbox{\parbox{0.9\textwidth}{\centering\vspace{2cm}
[Insert data flow diagram here]\\
Show: Source → Ingestion → Storage → Transformation → Analytics → Visualization
\vspace{2cm}}}
\caption{Data flow in the platform.}
\label{fig:data_flow}
\end{figure}

The data flow within the proposed platform embodies a cohesive pipeline that transforms raw, heterogeneous data into structured, semantically rich insights. It begins at the Data Source Layer, where data is ingested from multiple origins. The data then traverses the Data Ingestion Layer, which standardizes and streams the inputs. Upon ingestion, data is persisted in the Data Storage Layer. The Data Transformation Layer applies cleansing, normalization, and enrichment logic. Finally, insights propagate to the Analytics Layer where dashboards and reports are generated.

% ============================================================
% SECTION 2.3: TECHNOLOGY ARCHITECTURE
% ============================================================
\section{Technology Architecture}
\label{sec:tech_architecture}

% ==========================================================
% SUBSECTION 2.3.1: Data Source Technologies
% ==========================================================
\subsection{Data Source Technologies}
\label{sec:tech_source}

% PLACEHOLDER: MySQL, SQL Server
% Tham khảo thesis 2023/2026: mô tả chi tiết 2 DBMS

\textbf{MySQL:} MySQL is an open-source database management system. It is widely used due to its ease of deployment, speed, and compatibility across various systems. MySQL supports a wide range of storage engines and replication techniques, making it suitable for distributed environments and data-driven applications.

\textbf{Microsoft SQL Server:} Microsoft SQL Server is a relational database engine designed to retrieve, store, and manage data for structured querying, transactional processing and business intelligence features. SQL Server integrates seamlessly with enterprise tools such as Power BI and Azure Synapse.

% ==========================================================
% SUBSECTION 2.3.2: Data Ingestion Technologies
% ==========================================================
\subsection{Data Ingestion Technologies}
\label{sec:tech_ingestion}

% PLACEHOLDER: Debezium + Kafka / Airbyte
% Tham khảo thesis 2023: Airbyte
% Tham khảo thesis 2026: Debezium + Kafka

\textbf{[Tool A, e.g., Debezium or Airbyte]:} [Tool] is an open-source [type of tool] that helps [what it does]. One of the key advantages of [tool] is its capacity to manage a wide range of data sources and formats, including databases, APIs, and flat files.

\Cref{fig:ingestion_arch} shows the architecture of [tool].

\begin{figure}[tbp]
\centering
\fbox{\parbox{0.7\textwidth}{\centering\vspace{2cm}
[Insert tool architecture diagram here]
\vspace{2cm}}}
\caption{[Tool] architecture.}
\label{fig:ingestion_arch}
\end{figure}

\textbf{Apache Kafka:} Apache Kafka is a distributed event streaming platform that serves as the backbone in real-time data pipelines and streaming applications. Kafka is increasingly adopted in real-time analytics applications and event-driven architectures due to its ability to handle large volumes of data with high throughput.

Key components of Kafka include:
\begin{itemize}
    \item \textbf{Producers:} Clients that publish data to Kafka topics.
    \item \textbf{Topics:} Logical channels to which records are written.
    \item \textbf{Brokers:} Kafka servers that handle incoming data and client requests.
    \item \textbf{Consumers:} Applications that subscribe to topics and process messages.
\end{itemize}

% ==========================================================
% SUBSECTION 2.3.3: Storage and Lakehouse Technologies
% ==========================================================
\subsection{Storage and Lakehouse Technologies}
\label{sec:tech_storage}

% PLACEHOLDER: MinIO + Delta Lake / Apache Hudi

\textbf{MinIO:} MinIO is an open-source object storage system ideal for data storage applications such as data lakes and lakehouses. MinIO's distributed architecture enables it to scale horizontally, adding storage space and throughput as required.

\textbf{[Format Framework, e.g., Delta Lake or Apache Hudi]:} [Format] is a [description]. It brings database-like functionalities to large-scale data lakes, introducing capabilities such as record-level insert/update/delete (upsert), change data capture (CDC), and time travel queries.

% ==========================================================
% SUBSECTION 2.3.4: Data Transformation and Orchestration
% ==========================================================
\subsection{Data Transformation and Orchestration}
\label{sec:tech_transformation}

% PLACEHOLDER: Apache Spark + Apache Airflow

\textbf{Apache Spark:} Apache Spark is a widely used data processing engine particularly well-suited for large-scale data transformation use cases. Spark provides a fast and flexible platform for processing data in parallel across distributed clusters.

Key components of Spark:
\begin{itemize}
    \item \textbf{Driver Program:} The central coordinator that translates user code into tasks.
    \item \textbf{Cluster Manager:} Allocates resources across Spark applications.
    \item \textbf{Executors:} Distributed agents responsible for executing tasks.
\end{itemize}

\textbf{Apache Airflow:} Apache Airflow is an open-source platform for developing, scheduling, and monitoring batch workflows. It provides a Python-based framework where complex processes are expressed as Directed Acyclic Graphs (DAGs).

% ==========================================================
% SUBSECTION 2.3.5: Metadata and Knowledge Graph Management
% ==========================================================
\subsection{Metadata and Knowledge Graph Management}
\label{sec:tech_metadata}

% PLACEHOLDER: Datahub + Great Expectations

\textbf{DataHub:} DataHub is a modern data catalog designed to streamline metadata management, data discovery, and data governance. It enables efficient exploration of organizational data, lineage tracking, and enforcement of data contracts. DataHub models metadata as a graph of interconnected entities (datasets, dashboards, users) and relationships (schema, ownership, lineage).

\textbf{Great Expectations:} Great Expectations is an open-source data validation, documentation, and profiling tool designed to help teams maintain data quality throughout the data pipeline. It offers a declarative framework where users define expectations---assertions about data characteristics---that are automatically checked during data ingestion, transformation, and serving.

% ==========================================================
% SUBSECTION 2.3.6: Querying and Analytics
% ==========================================================
\subsection{Querying and Analytics}
\label{sec:tech_querying}

% PLACEHOLDER: Trino + Metabase / Superset

\textbf{Trino:} Trino, formerly known as PrestoSQL, is an open-source distributed SQL query engine designed for running fast, interactive analytic queries across heterogeneous data sources. It allows querying data directly where it lives---whether on object storage systems like S3, distributed file systems like HDFS, or in traditional relational databases---without requiring prior data movement.

The Trino architecture follows a distributed design based on a coordinator-worker model:
\begin{itemize}
    \item \textbf{Coordinator:} Parses SQL statements, optimizes query plans, schedules tasks.
    \item \textbf{Workers:} Execute parts of the query in parallel.
    \item \textbf{Catalogs:} Provide mappings to external data sources.
\end{itemize}

\textbf{[BI Tool, e.g., Metabase or Superset]:} [Tool] is an open-source business intelligence and data visualization platform designed to make data exploration accessible to non-technical users through an intuitive, low-code interface.

% ==========================================================
% SUBSECTION 2.3.7: LLM-Powered Analytics (Optional)
% ==========================================================
\subsection{LLM-Powered Intelligent Analytics (Optional)}
\label{sec:tech_llm}

The emergence of Large Language Models (LLMs) has fundamentally transformed the way users interact with data systems, introducing natural language interfaces that enable intuitive and dynamic data exploration.

In the context of intelligent analytics platforms, LLMs play a pivotal role by interpreting user prompts, automatically generating SQL queries, recommending visualizations, and providing narrative summaries of analytical results.

The architecture for LLM-powered intelligent analytics integrates LLMs into the data platform via a specialized interaction layer:
\begin{itemize}
    \item \textbf{Large Language Models:} Serve as the core natural language understanding and generation engine.
    \item \textbf{Integration Layer:} Acts as a bridge between the LLM and backend services.
\end{itemize}

% ============================================================
% SECTION 2.4: CHAPTER SUMMARY
% ============================================================
\section{Chapter Summary}
\label{sec:ch2_summary}

This chapter introduced the core layers and technology stack underpinning the proposed [Knowledge Graph-based / Data Fabric-based] Data Platform for [Intelligent Analytics / Modern Data Management]. We explored each component's role from data storage and transformation, to metadata management, intelligent querying, and LLM-powered analytics.

The proposed architecture consists of [5-7] major layers that work together to provide [capabilities]. The technology stack includes open-source tools such as [list of tools], which were selected based on their [criteria: scalability, flexibility, integration capability, community support].

Together, these technologies form an integrated ecosystem that enables efficient, scalable, and intelligent data workflows, providing a strong foundation for building advanced analytics capabilities within modern data-driven organizations. The following chapter will present the experimental setup and evaluation results.

\clearpage

% --- CHAPTER 3 ---
\chapter{Experiments and Evaluation}
\label{ch:experiments}
% =============================================================================
% CHAPTER 3: EXPERIMENTS AND EVALUATION
% =============================================================================
% COMPLETELY REWRITTEN — Old Knowledge Graph / Data Fabric content removed
% Project: A Context-Aware Framework for Streaming Data Quality Monitoring
% Date: April 24, 2026
% =============================================================================

\section{Introduction to Chapter}
\label{sec:ch3-intro}

This chapter describes the experimental evaluation of StreamDQ. We begin with the
experimental setup (\ref{sec:experimental-setup}): datasets, infrastructure,
and synthetic anomaly injection. \ref{sec:evaluation-methodology} presents the
ground-truth tracking protocol, metrics, ablation design, and statistical tests.
\ref{sec:results-syn-sem} through \ref{sec:results-latency} present benchmark
results as placeholder tables labeled \texttt{[TIER-2\ ESTIMATED]}, pending
implementation and benchmark execution. \ref{sec:ch3-limitations} presents
the mandatory limitations section with all seven documented constraints.

\section{Experimental Setup}
\label{sec:experimental-setup}

\subsection{Datasets}

\textbf{NYC TLC Yellow Taxi.} We use the NYC Taxi and Limousine Commission
Yellow Taxi trip record dataset for January 2024 \citep{nyc_tlc}, stored as
Parquet files with approximately 3 million records. These records contain
zone-level spatial references (PULocationID and DOLocationID, mapping to 263
taxi zones) but \textbf{no GPS latitude/longitude coordinates}. The TLC
dataset is replayed through Apache Kafka at real-time speed using a custom
Kafka producer. Because it lacks GPS, only SYN and SEM rules apply; CRS rules
cannot be evaluated on TLC data.

\textbf{NYC MTA Bus GTFS-realtime.} We use the live NYC MTA Bus GTFS-realtime
public feed \citep{nyc_mta_bus}, which publishes VehiclePosition messages for
approximately 500 buses at approximately 30-second intervals per vehicle.
This feed is publicly accessible and requires no API key. The feed URL,
evaluation timestamp, and dataset version are logged for reproducibility.
Because the MTA Bus feed includes GPS latitude/longitude coordinates, all
three rule layers (SYN, SEM, and CRS) apply to this stream.

Figure \ref{fig:experimental-pipeline} shows the evaluation pipeline.

\begin{figure}[htbp]
  \centering
  \fbox{\parbox{0.92\linewidth}{
    \centering
    \textbf{[Figure F7: Experimental Pipeline — see
    \texttt{figures/F7\_Experimental\_Pipeline.excalidraw.md}]}
  }}
  \caption{Evaluation pipeline for StreamDQ. Synthetic anomalies are injected
    at controlled rates into the event stream before rule evaluation. Every
    injected anomaly is tracked in the ground\_truth\_events table. Every
    detected violation is matched against ground truth to compute precision
    and recall.}
  \label{fig:experimental-pipeline}
\end{figure}

\subsection{Infrastructure}

StreamDQ operates in two execution modes:

\textbf{LocalPipeline (development mode).} An in-process Python pipeline
processes events from Kafka or a replay log. State is maintained in Python
dictionaries (in-memory). SQLite is used for violation storage. LocalPipeline
is used for unit tests, development, and functional validation.
\texttt{[TIER-2\ ESTIMATED]}

\textbf{FlinkPipeline (distributed mode).} Apache Flink processes events
from Kafka in a distributed cluster. State is maintained in RocksDB.
PostgreSQL stores violations and ground-truth data. FlinkPipeline is used
for full benchmark runs.
\texttt{[TIER-2\ ESTIMATED]}

\textbf{Critical distinction.} LocalPipeline and FlinkPipeline have
different latency and throughput characteristics. LocalPipeline results are
not equivalent to distributed Flink results. All latency numbers reported
from LocalPipeline are labeled ``LocalPipeline profiling only, not
equivalent to distributed Flink performance.''

Figure \ref{fig:dashboard-mockup} shows the monitoring dashboard.

\begin{figure}[htbp]
  \centering
  \fbox{\parbox{0.92\linewidth}{
    \centering
    \textbf{[Figure F8: Grafana Dashboard Mockup — see
    \texttt{figures/F8\_Dashboard\_Mockup.excalidraw.md}]}
  }}
  \caption{Grafana dashboard for StreamDQ monitoring. Four panels display:
    (1) violation rate by rule, (2) TQS composite score over time,
    (3) L0 coverage fraction, and (4) processing latency percentiles.
    Screenshot placeholder: replace with actual dashboard after deployment.}
  \label{fig:dashboard-mockup}
\end{figure}

\subsection{Synthetic Anomaly Injection}

We inject synthetic anomalies into the event stream to create ground truth
for evaluation. Each injected anomaly is tagged with its entity index,
anomaly type, and injection rate. Table \ref{tab:anomaly-types} summarizes
the anomaly types and injection parameters.

\begin{table}[htbp]
  \centering
  \caption{Synthetic anomaly types and injection parameters.
    NG-4: CRS003 recall is \texttt{[TIER-3\ UNMEASURABLE]}
    because the replay suppression gate blocks synthetic duplicate injection.}
  \label{tab:anomaly-types}
  \begin{tabular}{lllr}
    \toprule
    Anomaly Type & Rule Triggered & Injection Method & Status \\
    \midrule
    NULL field         & SYN001 & Set required field to NULL       & Ready \\
    NaN field          & SYN001 & Set float field to NaN          & Ready \\
    Negative fare      & SYN002 & Set fare\_amount $<$ 0           & Ready \\
    GPS speed spike   & CRS001 & Two positions 2 km apart in 30 s & Ready \\
    GPS jump          & CRS002 & Positions $>$400 m at $\leq$30 s & Ready \\
    Duplicate event   & CRS003 & Emit original + duplicate         & NG-4 blocked \\
    \bottomrule
  \end{tabular}
\end{table}

Injection rates are 0.5\%, 1.0\%, 2.0\%, and 5.0\% of events per anomaly type.
Each trial uses a different random seed (SEED = 42, 43, 44). A warmup period
of 10,000 events (approximately 60 seconds at real-time speed) precedes each
measurement window. The measurement window is 600 seconds (10 minutes) per trial.
Three independent trials are run per condition.

\textbf{NG-4 (CRS003 blocker).} The replay suppression gate in the current
code tags synthetic duplicate events as \texttt{is\_replay=True}, which triggers
the gate in the deduplication logic and suppresses violation emission. To measure
CRS003 recall, synthetic duplicate events must be tagged \texttt{is\_replay=False}.
This is an evaluation infrastructure bug, not a rule logic bug. CRS003 recall
is labeled \texttt{[TIER-3\ UNMEASURABLE]} until NG-4 is fixed.

\section{Evaluation Methodology}
\label{sec:evaluation-methodology}

\subsection{Ground-Truth Tracking}

We track every injected anomaly in a PostgreSQL \texttt{ground\_truth\_events}
table with the following schema:

\begin{itemize}
  \itemsep0em
  \item \texttt{run\_id}: UUID identifying the evaluation run
  \item \texttt{injection\_rate}: float (0.005, 0.01, 0.02, 0.05)
  \item \texttt{anomaly\_type}: VARCHAR (NULL, NAN, NEGATIVE\_FARE, GPS\_SPEED,
        GPS\_JUMP, DUPLICATE)
  \item \texttt{entity\_index}: INTEGER (position in event stream)
  \item \texttt{injected\_value}: JSONB (field name, injected value)
  \item \texttt{injected\_at}: TIMESTAMP
\end{itemize}

Every detected violation is matched against ground truth using exact
entity index matching. No fuzzy matching is used.

\subsection{Metrics}

Table \ref{tab:eval-metrics} presents the evaluation metrics.

\begin{table}[htbp]
  \centering
  \caption{Evaluation metrics, confidence interval methods, and bootstrap
    specifications. Bootstrap iterations should be increased to 10,000 for
    publication.}
  \label{tab:eval-metrics}
  \begin{tabular}{lllr}
    \toprule
    Metric & CI Method & Bootstrap Iterations & Formula \\
    \midrule
    Precision, Recall, F1 & Wilson score     & $\geq$ 1,000 & $2P\!R/(P+R)$ \\
    Violation rate        & Clopper-Pearson & $\geq$ 1,000 & $n_{\text{viol}}/N$ \\
    P99 latency           & Percentile       & $\geq$ 1,000 & $P_{99}(L)$ \\
    Throughput            & t-test bootstrap & $\geq$ 1,000 & events/s \\
    TQS composite         & Percentile       & $\geq$ 1,000 & Weighted sum \\
    Spearman $\rho$       & Fisher z-transform& $\geq$ 1,000 & Rank corr. \\
    \bottomrule
  \end{tabular}
\end{table}

All confidence intervals are at the 95\% level using the percentile bootstrap
method (1,000 iterations). For publication, bootstrap iterations should be
increased to 10,000.
\texttt{[TODO: increase bootstrap to 10,000 iterations for publication]}

\textbf{Latency measurement.} Processing latency is measured from event arrival
(at Kafka produce timestamp) to violation stored (PostgreSQL INSERT complete).
This is an end-to-end measurement, not an internal processing time.
Latency numbers from LocalPipeline are labeled as such.

\textbf{TQS correlation.} TQS composite scores are correlated with downstream
ETA prediction error, not with injection rate. ETA prediction error is
computed independently from GTFS static schedule data and published real-time
delays from the MTA API. This design avoids the circularity of measuring
TQS against injection rate.

\subsection{Ablation Studies}

We conduct three ablation studies to answer the research questions:

\begin{itemize}
  \itemsep0em
  \item \textbf{RQ1 ablation}: Static (global L4) thresholds vs. L0--L4
        context-aware thresholds. Measures whether finer context improves
        detection F1 at L0 cells.
  \item \textbf{RQ4 ablation}: SYN-only vs. SYN+SEM vs. SYN+SEM+CRS.
        Measures whether each rule layer adds incremental F1.
  \item \textbf{RQ6 ablation}: Rule-only vs. rule+ML (Bayesian Optimization only).
        Measures whether ML calibration improves F1.
\end{itemize}

\subsection{Statistical Tests}

Table \ref{tab:statistical-tests} presents the statistical tests.

\begin{table}[htbp]
  \centering
  \caption{Statistical tests, significance levels, and corrections.
    Holm--Bonferroni is applied within test families.}
  \label{tab:statistical-tests}
  \begin{tabular}{lllr}
    \toprule
    Test & Purpose & $\alpha$ & Correction \\
    \midrule
    Wilcoxon signed-rank & Non-parametric paired comparison & 0.05 & Holm--Bonferroni \\
    McNemar's mid-p exact & Paired binary classifier comparison & 0.01 & Holm--Bonferroni \\
    Friedman's test & Multiple related samples (L0--L4) & 0.05 & Nemenyi post-hoc \\
    Bootstrap percentile & Non-parametric CI & 0.05 & Per-method \\
    \bottomrule
  \end{tabular}
\end{table}

Holm--Bonferroni correction is applied within test families at $\alpha = 0.05$.
For RQ1 (5 pairwise comparisons), $\alpha_{\text{adj}} = 0.05 / 5 = 0.01$.
For RQ2 (6 pairwise comparisons), $\alpha_{\text{adj}} = 0.05 / 6 \approx 0.0083$.

\section{Results: SYN/SEM Rules}
\label{sec:results-syn-sem}

Table \ref{tab:results-syn-sem} presents precision, recall, and F1 for SYN
and SEM rules on the NYC TLC Yellow Taxi dataset.

\begin{table}[htbp]
  \centering
  \caption{Precision, recall, and F1 for SYN and SEM rules on NYC TLC.
    Results are \texttt{[TIER-2\ ESTIMATED]} pending benchmark execution.
    All values are accompanied by 95\% bootstrap confidence intervals
    (1,000 iterations). TODO: increase bootstrap to 10,000 iterations.}
  \label{tab:results-syn-sem}
  \begin{tabular}{lcccccc}
    \toprule
    \multirow{2}{*}{Rule} & \multirow{2}{*}{Precision} & \multirow{2}{*}{Recall}
    & \multirow{2}{*}{F1} & \multirow{2}{*}{95\% CI} &
    \multirow{2}{*}{n$_{\text{inj}}$} \\
    & & & & & \\
    \midrule
    SYN001 (NULL)      & [TIER-2\ ESTIMATED] & [TIER-2\ ESTIMATED]
                       & [TIER-2\ ESTIMATED] & — & — \\[4pt]
    SYN001 (NaN)       & [TIER-2\ ESTIMATED] & [TIER-2\ ESTIMATED]
                       & [TIER-2\ ESTIMATED] & — & — \\[4pt]
    SYN002 (neg fare)   & [TIER-2\ ESTIMATED] & [TIER-2\ ESTIMATED]
                       & [TIER-2\ ESTIMATED] & — & — \\[4pt]
    SYN003 (bounds)    & [TIER-2\ ESTIMATED] & [TIER-2\ ESTIMATED]
                       & [TIER-2\ ESTIMATED] & — & — \\[4pt]
    SEM001 (fare)      & [TIER-2\ ESTIMATED] & [TIER-2\ ESTIMATED]
                       & [TIER-2\ ESTIMATED] & — & — \\[4pt]
    SEM002 (distance)  & [TIER-2\ ESTIMATED] & [TIER-2\ ESTIMATED]
                       & [TIER-2\ ESTIMATED] & — & — \\[4pt]
    SEM003 (pax count) & [TIER-2\ ESTIMATED] & [TIER-2\ ESTIMATED]
                       & [TIER-2\ ESTIMATED] & — & — \\[4pt]
    GTFSSem002 (stale) & [TIER-2\ ESTIMATED] & [TIER-2\ ESTIMATED]
                       & [TIER-2\ ESTIMATED] & — & — \\
    \bottomrule
  \end{tabular}
\end{table}

The SYN001 (NULL) and SYN001 (NaN) rules provide the baseline precision
for record-level validation. We expect high precision (close to 1.0) for
NULL checks because synthetic NULL injection produces unambiguous violations.
SEM001 and SEM002 are expected to have lower recall because context-aware
thresholds may not fire on injected anomalies that fall within the rolling
percentile bounds.

\section{Results: CRS Rules}
\label{sec:results-crs}

Table \ref{tab:results-crs} presents precision, recall, and F1 for CRS
rules on the NYC MTA Bus GTFS-realtime stream.

\begin{table}[htbp]
  \centering
  \caption{Precision, recall, and F1 for CRS rules on NYC MTA Bus GTFS-realtime.
    CRS001 and CRS002 are \texttt{[TIER-2\ ESTIMATED]} pending benchmark.
    CRS003 is \texttt{[TIER-3\ UNMEASURABLE]} due to NG-4
    (replay suppression gate blocks synthetic duplicate injection).
    All values are accompanied by 95\% bootstrap confidence intervals
    (1,000 iterations). TODO: increase bootstrap to 10,000 iterations.}
  \label{tab:results-crs}
  \begin{tabular}{lcccccc}
    \toprule
    \multirow{2}{*}{Rule} & \multirow{2}{*}{Precision} & \multirow{2}{*}{Recall}
    & \multirow{2}{*}{F1} & \multirow{2}{*}{95\% CI} &
    \multirow{2}{*}{n$_{\text{inj}}$} \\
    & & & & & \\
    \midrule
    CRS001 (GPS speed)   & [TIER-2\ ESTIMATED] & [TIER-2\ ESTIMATED]
                         & [TIER-2\ ESTIMATED] & — & — \\[4pt]
    CRS002 (GPS jump)    & [TIER-2\ ESTIMATED] & [TIER-2\ ESTIMATED]
                         & [TIER-2\ ESTIMATED] & — & — \\[4pt]
    CRS003 (deduplication)& [TIER-2\ ESTIMATED] & \cellcolor[gray]{0.85}
                         \texttt{[TIER-3]} & \cellcolor[gray]{0.85}
                         \texttt{[TIER-3]} & \cellcolor[gray]{0.85} NG-4
                         & \cellcolor[gray]{0.85} — \\
    \bottomrule
  \end{tabular}
\end{table}

CRS001 and CRS002 use physics-based thresholds (speed bounds, jump distance)
that are independent of the data distribution, so we expect high precision.
CRS001 lower bound (2.0 km/h) may produce false positives for stationary
vehicles with GPS jitter; the sustained-low-speed guard ($\Delta t > 60$ s)
mitigates this. CRS002 precision depends on the quality of the synthetic
GPS jump injection; the 400-meter threshold was calibrated for the 30-second
GTFS-realtime update interval.

CRS003 recall is \texttt{[TIER-3\ UNMEASURABLE]} because the replay
suppression gate (NG-4) blocks synthetic duplicate injection. Fixing NG-4
requires tagging synthetic duplicate events as \texttt{is\_replay=False}.

\section{Results: Context-Aware Threshold Ablation}
\label{sec:results-threshold-abl}

Table \ref{tab:results-threshold-abl} presents the ablation results for
context-aware thresholds (RQ1).

\begin{table}[htbp]
  \centering
  \caption{Ablation: static (L4 global) vs. context-aware (L0--L4) thresholds.
    ``$\Delta$F1 at L0 cells'' measures the F1 improvement specifically for
    events falling into L0 context cells.
    ``Aggregate $\Delta$F1'' measures improvement across all events.
    L0 coverage is 2--5\%; aggregate F1 improvement is expected at 1--2 pp.
    Results are \texttt{[TIER-2\ ESTIMATED]} pending benchmark execution.}
  \label{tab:results-threshold-abl}
  \begin{tabular}{lcccccc}
    \toprule
    Configuration & $\Delta$F1 at L0 cells (pp) & 95\% CI
    & Aggregate $\Delta$F1 (pp) & 95\% CI & L0 coverage \\
    \midrule
    Static (L4 global)   & baseline & — & baseline & — & — \\
    L0 only              & [TIER-2] & — & [TIER-2] & — & 2--5\% \\
    L0--L4 (hierarchical)& [TIER-2] & — & [TIER-2] & — & 2--5\% \\
    \bottomrule
  \end{tabular}
\end{table}

\textbf{Critical caveat.} The 5 pp target for $\Delta$F1 applies specifically
to events falling into L0 context cells (2--5\% of events). Because most
events fall to L3 or L4 (where thresholds are similar to static), the
aggregate $\Delta$F1 across all events is expected to be 1--2 pp, not 5 pp.
Both metrics will be reported separately. The aggregate number is likely
underpowered at the current sample size (approximately 90 measurement windows
$\times$ 3 trials).

Statistical significance is assessed via Wilcoxon signed-rank test with
Bonferroni correction ($\alpha_{\text{adj}} = 0.01$ for 5 comparisons).

\section{Results: ML Augmentation}
\label{sec:results-ml}

Table \ref{tab:results-ml} presents the ablation results for ML-augmented
threshold calibration (RQ6).

\begin{table}[htbp]
  \centering
  \caption{Ablation: rule-only vs. rule+Bayesian Optimization.
    Bayesian Optimization runs hourly, tuning k\_multiplier, IF\_alpha,
    and weekend\_discount on F1 objective (not violation rate).
    Isolation Forest is conditional (GO if $\rho_{\text{IF,P90}} < 0.8$).
    XGBoost is scaffold only (training target undefined).
    LSTM is NO-GO (GPS archive unconfirmed; high CRS002 redundancy).
    Results are \texttt{[TIER-2\ ESTIMATED]} pending benchmark execution.
    Positive, negative, and mixed results are all valid findings.}
  \label{tab:results-ml}
  \begin{tabular}{lcccccc}
    \toprule
    Configuration & Precision & Recall & F1 & 95\% CI & n$_{\text{inj}}$ \\
    \midrule
    Rule-only (baseline) & [TIER-2] & [TIER-2] & [TIER-2] & — & — \\
    + Bayesian Optimization & [TIER-2] & [TIER-2] & [TIER-2] & — & — \\
    + Isolation Forest (cond.) & [TIER-2] & [TIER-2] & [TIER-2] & — & — \\
    \bottomrule
  \end{tabular}
\end{table}

This is a \textit{measured contribution}: if ML provides no improvement,
the negative result is scientifically valuable and will be reported as such.

\section{Latency and Throughput}
\label{sec:results-latency}

Table \ref{tab:results-latency} presents latency and throughput measurements.

\begin{table}[htbp]
  \centering
  \caption{Latency and throughput measurements. LocalPipeline results are
    \texttt{[TIER-2\ ESTIMATED]} and labeled ``LocalPipeline profiling only,
    not equivalent to distributed Flink performance.'' Distributed Flink
    benchmarks require infrastructure deployment.}
  \label{tab:results-latency}
  \begin{tabular}{lcccc}
    \toprule
    Mode & P50 latency (ms) & P99 latency (ms) & 95\% CI (ms) &
    Throughput (events/s) \\
    \midrule
    LocalPipeline & [TIER-2] & [TIER-2] & — & [TIER-2] \\
    FlinkPipeline & [TIER-2] & [TIER-2] & — & [TIER-2] \\
    \bottomrule
  \end{tabular}
\end{table}

\textbf{Critical limitation.} LocalPipeline and FlinkPipeline have fundamentally
different architectures. LocalPipeline processes events in-process with Python
state management; FlinkPipeline distributes events across JVM TaskManagers with
RocksDB state. P99 latency from LocalPipeline is not representative of
distributed Flink latency. Distributed benchmarks require Flink cluster
deployment and are not yet available.

\section{Limitations}
\label{sec:ch3-limitations}

We explicitly acknowledge seven limitations of this work.

\textbf{1. CRS003 recall is unmeasurable (NG-4).} The replay suppression
gate tags synthetic duplicate events as \texttt{is\_replay=True}, which
triggers the gate and suppresses violation emission. CRS003 recall
cannot be measured until synthetic duplicates are tagged
\texttt{is\_replay=False}. This is an evaluation infrastructure bug, not
a rule logic bug. CRS003 precision is measurable; recall is labeled
\texttt{[TIER-3\ UNMEASURABLE]}.
\texttt{[TIER-3\ UNMEASURABLE]}

\textbf{2. L0 coverage is sparse (2--5\% of events).} The aggregate
$\Delta$F1 from context-aware thresholds across all events is likely
1--2 pp, not the 5 pp target which applies only to L0 cells. Both
L0-specific and aggregate metrics will be reported separately. The aggregate
metric is likely underpowered at approximately 90 measurement windows.

\textbf{3. ML augmentation results are estimated (TIER-2).} All ML
numbers (F1 improvement from Bayesian Optimization, Isolation Forest
conditional activation) require Phase 3 implementation and benchmark
execution. These results are placeholders.

\textbf{4. Single dataset (NYC TLC + NYC MTA Bus).} Results may not
generalize to other transit agencies, geographic regions, or data schemas.
CRS rules (CRS001, CRS002) calibrated for NYC MTA Bus may need parameter
adjustment for different vehicle types, update frequencies, or geographic
contexts.

\textbf{5. LocalPipeline $\neq$ distributed Flink.} All latency and
throughput numbers from LocalPipeline are labeled as development-mode
profiling only. Distributed Flink benchmarks require cluster deployment
and are not yet available. Claiming Flink-level latency from LocalPipeline
results would be architecturally incorrect.

\textbf{6. L4 global fallback masks context-specific patterns.} Events
evaluated at L4 use the global threshold regardless of their context.
If the global threshold is lenient, L4 events that would violate a
context-specific threshold will not be flagged. This is an inherent
trade-off of the hierarchical fallback design: it ensures coverage but
dilutes context specificity.

\textbf{7. D4 External context is a stub.} The holiday indicator in
D4 (External context dimension) is not implemented. The five-dimensional
context decomposition effectively operates on four dimensions (D1, D2, D3,
D5). Holidays, traffic incidents, and other external signals are not
captured by the current context model.

\section{Chapter Summary}
\label{sec:ch3-summary}

This chapter presented the evaluation methodology and results. We described
the NYC TLC and NYC MTA Bus datasets, the LocalPipeline and FlinkPipeline
infrastructure, and the synthetic anomaly injection protocol with five anomaly
types. We detailed ground-truth tracking, metrics with 95\% bootstrap CI,
ablation studies for RQ1, RQ4, and RQ6, and statistical tests with
Holm--Bonferroni correction. Result tables are placeholders labeled
\texttt{[TIER-2\ ESTIMATED]} pending benchmark execution; CRS003 is
\texttt{[TIER-3\ UNMEASURABLE]} due to NG-4. The mandatory limitations
section explicitly acknowledged all seven documented constraints.

\clearpage

% ============================================================
% CONCLUSION AND PERSPECTIVES
% ============================================================
\chapter*{Conclusion and Perspectives}
\label{ch:conclusion}
\addcontentsline{toc}{chapter}{Conclusion and Perspectives}
% =============================================================================
% CONCLUSION AND PERSPECTIVES
% =============================================================================
% COMPLETELY REWRITTEN — Old Knowledge Graph / Data Fabric content removed
% Project: A Context-Aware Framework for Streaming Data Quality Monitoring
% Date: April 24, 2026
% =============================================================================

\section{Summary of Contributions}
\label{sec:conclusion-contributions}

We presented StreamDQ: a context-aware framework for streaming data quality
monitoring, implemented in Apache Flink and evaluated on NYC transportation
data streams. We summarize the four primary contributions of this work.

\textbf{C1: First streaming GPS trajectory validation on real GTFS-realtime
feeds.} We implemented three cross-record rules — CRS001 (Haversine-based GPS
speed bounds $[2.0,\ 100.0]$ km/h), CRS002 (GPS jump detection $>400$ m
in $\leq 30$ s), and CRS003 (event deduplication with 300-second window) —
as Java \texttt{KeyedProcessFunction} operators invoked via gRPC from the
Flink Python pipeline. These rules are evaluated exclusively on the NYC MTA
Bus GTFS-realtime stream, which provides GPS coordinates; they cannot be
evaluated on the NYC TLC Yellow Taxi dataset (zone-level only).
\texttt{[TIER-2\ ESTIMATED]}

\textbf{C2: Hierarchical context-aware threshold fallback (L0--L5).} We
designed a five-level hierarchical threshold system that traverses from the
finest-grained context cell (L0: hour $\times$ zone $\times$ weekday/weekend)
to physics priors (L5), with graceful degradation through L1--L4. The
minimum sample thresholds per level (L0 $\geq 100$, L1 $\geq 50$, L2 $\geq 25$,
L3 $\geq 10$, L4 $\geq 5$) are derived from power analysis (Cohen's $d = 0.30$,
$\alpha = 0.01$, power $= 0.80$). We validated that this hierarchical
fallback mechanism is absent from all surveyed frameworks (Great Expectations,
Soda Core, Stream DaQ, METER, Ada-Context, Deequ, GTFS Validator, and
7 others). Expected L0 coverage is 2--5\% of events due to zone-hour sparsity
in the NYC TLC dataset.
\texttt{[TIER-2\ ESTIMATED]}

\textbf{C3: Ground-truth evaluation methodology.} We implemented a complete
evaluation pipeline comprising synthetic anomaly injection, per-entity ground-truth
tracking in PostgreSQL, and precision/recall/F1 metrics with 95\% bootstrap
confidence intervals (1,000 iterations). We explicitly acknowledge what cannot
be measured: CRS003 recall is blocked by NG-4 (replay suppression gate), and
distributed Flink latency is not benchmarked. This honest methodology is
modeled on Exathlon \citep{exathlon2021}, which established precedent for
benchmarking streaming anomaly detection.
\texttt{[TIER-1\ VERIFIED\ for\ methodology;\ TIER-2\ ESTIMATED\ for\ results]}

\textbf{C4: ML-augmented threshold calibration.} We designed a three-component
ML pipeline — Bayesian Optimization (Priority 1, GO), Isolation Forest (Priority 2,
conditional), and XGBoost (Priority 3, conditional) — that calibrates context
thresholds without overriding rule-based violation detection. Rules remain
authoritative at all times. LSTM (Priority 4, NO-GO) is excluded due to
unconfirmed GPS training data and high redundancy with CRS002.
\texttt{[TIER-2\ ESTIMATED]}

\section{Key Findings}
\label{sec:conclusion-findings}

We report the following key findings based on the design and methodology
established in this thesis.

\textbf{CRS001 and CRS002 are designed as specified with physics-based
thresholds.} CRS001 uses Haversine distance with the $[2.0,\ 100.0]$ km/h
bound calibrated for NYC MTA Bus urban operations. The lower bound of
2.0 km/h excludes GPS jitter; the upper bound of 100.0 km/h closes the
detection gap for moderate spoofing (20--100 km/h). CRS002 uses the
400-meter/30-second threshold calibrated for the 30-second GTFS-realtime
update interval. Both rules use L5 physics priors exclusively; they have
no context dependency.
\texttt{[TIER-2\ ESTIMATED]}

\textbf{Context-aware thresholds provide value at L0 cells but coverage is
sparse.} The expected L0 coverage of 2--5\% of events means that most
events are evaluated at coarser fallback levels (L3 or L4). At L4 (global
fallback), the threshold is identical to a static threshold, so no
context-specific improvement is available. The aggregate $\Delta$F1 across
all events is expected at 1--2 pp, not the 5 pp target which applies only
to L0 cells. This finding is critical for interpreting RQ1 honestly.
\texttt{[TIER-2\ ESTIMATED]}

\textbf{L4 global fallback is necessary but dilutes context specificity.}
The hierarchical fallback ensures that every event receives a threshold
(even if only the physics prior), preventing silent failures. However, the
L4 global threshold pools all events regardless of context, which may
increase the false-negative rate for events with unusual but legitimate
values. This trade-off is inherent to the fallback design.

\textbf{ML augmentation is a measured contribution.} Phase 3 implementation
is mandatory. Whether ML provides measurable F1 improvement over rule-only
thresholds is an empirical question that requires benchmark execution.
Positive, negative, and mixed results are all scientifically valid and
will be reported as such.

\textbf{TQS correlation with downstream quality remains to be validated.}
The non-circular TQS design (five stream-intrinsic dimensions, weighted
sum) distinguishes StreamDQ from prior work that defines TQS as
$1 - \text{violation\_rate}$. Whether TQS meaningfully predicts downstream
quality (e.g., ETA prediction accuracy) requires independent measurement
using ETA error computed from GTFS static schedules and published delays.
\texttt{[TIER-2\ ESTIMATED]}

\section{Future Work}
\label{sec:conclusion-future}

Five directions for future work emerge from this research.

\textbf{1. Fix NG-4: replay suppression gate for CRS003 evaluation.}
The current code tags synthetic duplicate events as \texttt{is\_replay=True},
which triggers the replay suppression gate and prevents violation emission.
Fixing this bug requires tagging synthetic duplicates as \texttt{is\_replay=False}.
Once fixed, CRS003 recall can be measured alongside precision, completing
the CRS rule evaluation.

\textbf{2. Implement D4: External context (holiday indicator).} The
five-dimensional context decomposition currently operates on four dimensions
(D1, D2, D3, D5). Implementing the holiday indicator in D4 would upgrade
the context model to full five-dimensional operation. Holidays affect
both taxi fare distributions (surcharges, airport trips) and bus ridership
patterns. A holiday lookup table for NYC and integration with the context
extraction pipeline are the primary implementation steps.

\textbf{3. Deploy Flink cluster for distributed benchmarks.} LocalPipeline
results are not equivalent to distributed Flink performance. Deploying a
multi-node Flink cluster with Kafka partitioning and RocksDB state would
enable proper latency and throughput measurements. This is necessary for
claiming P99 latency numbers that are representative of real-world
streaming deployment.

\textbf{4. Re-evaluate LSTM if GPS training data becomes available.}
The LSTM trajectory prediction component is explicitly excluded (Priority 4,
NO-GO) due to unconfirmed 6-month historical NYC MTA Bus GPS archive.
If this archive is confirmed to exist, LSTM should be re-evaluated under
two conditions: (a) CRS002 recall falls below 60\% on real GPS anomalies,
and (b) LSTM provides measurable improvement over CRS002 alone. LSTM
would serve only as an alert elevation mechanism, never as authoritative
validation.

\textbf{5. Extend CRS rules to other GTFS feeds beyond NYC.} The
CRS001 and CRS002 rules are calibrated for NYC MTA Bus operations (urban
bus, 30-second update interval, 100 km/h speed limit). Generalizing to
other transit agencies requires parameter recalibration for different
vehicle types (subway, rail, ferry), geographic contexts (rural vs. urban),
and update frequencies. A calibration framework that automates threshold
derivation from the GTFS-realtime feed characteristics would enable
broader deployment.

\section*{Chapter Summary}
\label{sec:conclusion-summary}

This thesis presented StreamDQ: a context-aware framework for streaming
data quality monitoring, designed for GPS trajectory validation on real
GTFS-realtime transit feeds. We implemented three layers of data quality
rules (syntactic, semantic, cross-record), a hierarchical context-aware
threshold system with L0--L5 fallback, a non-circular Trajectory Quality
Scoring mechanism, and ML-augmented threshold calibration. We designed a
rigorous evaluation methodology with synthetic anomaly injection, ground-truth
tracking, and bootstrap confidence intervals. We explicitly acknowledged
seven limitations, including the NG-4 blocker for CRS003 recall, the sparse
L0 coverage, and the distinction between LocalPipeline and distributed Flink
benchmarks. StreamDQ is a research and education platform that demonstrates
the feasibility and measurability of streaming GPS trajectory validation
for public transit data quality monitoring.

\clearpage

% ============================================================
% REFERENCES
% ============================================================
\backmatter
\bibliographystyle{plainnat}
\bibliography{references}
\clearpage

% ============================================================
% APPENDICES
% ============================================================
\begin{appendices}
\chapter{Additional Experimental Results}
\label{app:results}
% ============================================================
% APPENDIX A: Additional Experimental Results
% ============================================================

\chapter{Additional Experimental Results}
\label{app:results}

% PLACEHOLDER: Thêm các kết quả thực nghiệm bổ sung
% Ví dụ: bảng kết quả chi tiết hơn, đồ thị bổ sung, v.v.

\section{Extended Performance Metrics}
\label{app:extended_metrics}

\Cref{tab:extended_metrics} presents additional performance metrics from the experiments.

\begin{table}[tbp]
\centering
\caption{Extended performance metrics.}
\label{tab:extended_metrics}
\begin{tabular}{@{}lcccccc@{}}
\toprule
\textbf{Scenario} & \textbf{Records} & \textbf{Tables} & \textbf{Sync Time} & \textbf{CPU} & \textbf{Memory} \\
\midrule
Scenario 1 & 100,000 & 13 & $<1$ min & -- & -- \\
Scenario 2 & 2,000,000 & 13 & $\sim2$ min & -- & -- \\
Scenario 3 & 3,000 & 2 & $\sim10$ sec & -- & -- \\
\bottomrule
\end{tabular}
\end{table}

\section{Additional Data Lineage Diagrams}
\label{app:lineage_extra}

\Cref{fig:lineage_extra} shows an extended view of the data lineage graph.

\begin{figure}[tbp]
\centering
\fbox{\parbox{0.9\textwidth}{\centering\vspace{2cm}
[Insert additional data lineage screenshot here]
\vspace{2cm}}}
\caption{Extended data lineage view.}
\label{fig:lineage_extra}
\end{figure}


\chapter{Configuration Files}
\label{app:config}
% ============================================================
% APPENDIX B: Configuration Files
% ============================================================

\chapter{Configuration Files}
\label{app:config}

% PLACEHOLDER: Đặt các file cấu hình quan trọng
% Ví dụ: docker-compose.yml, Spark config, Airflow DAG, v.v.

\section{Docker Compose Configuration}
\label{app:docker_compose}

\Cref{fig:docker_compose} shows the Docker Compose configuration used for the platform.

\begin{figure}[tbp]
\centering
\fbox{\parbox{0.8\textwidth}{\centering\vspace{2cm}
[Insert docker-compose.yml content here]
\vspace{2cm}}}
\caption{Docker Compose configuration.}
\label{fig:docker_compose}
\end{figure}

\section{Spark Job Configuration}
\label{app:spark_config}

\Cref{fig:spark_config} shows the Spark job configuration.

\begin{figure}[tbp]
\centering
\fbox{\parbox{0.8\textwidth}{\centering\vspace{2cm}
[Insert Spark configuration or PySpark code here]
\vspace{2cm}}}
\caption{Spark job configuration.}
\label{fig:spark_config}
\end{figure}

\section{Airflow DAG Definition}
\label{app:airflow_dag}

\Cref{fig:airflow_dag_code} shows the Python code for the Airflow DAG.

\begin{figure}[tbp]
\centering
\fbox{\parbox{0.8\textwidth}{\centering\vspace{2cm}
[Insert Airflow DAG Python code here]
\vspace{2cm}}}
\caption{Airflow DAG definition.}
\label{fig:airflow_dag_code}
\end{figure}

\end{appendices}

\end{document}
